\documentclass{article}
\usepackage[T1]{fontenc}
\usepackage{lmodern}
\usepackage{graphicx}
\usepackage{amsmath,amssymb}
\usepackage[a4paper,margin=2cm]{geometry}
\usepackage[absolute,overlay]{textpos}
\setlength{\TPHorizModule}{1mm}
\setlength{\TPVertModule}{1mm}
\usepackage[indonesian]{babel}
\usepackage{hyperref}
\setlength{\emergencystretch}{2em}
% unit: Halaman 35

\begin{document}

% === Halaman 35 (python/native) ===

• Kunci adalah string: K = k1k2 … km

 ki untuk 1  i  m menyatakan huruf-huruf alfabet 

• Jika panjang kunci lebih pendek daripada panjang plainteks, maka kunci diulang  

secara periodik. 

• Misalkan panjang kunci m = 10, maka 10 huruf pertama plainteks dienkripsi 

dengan kunci K, setiap huruf ke-i menggunakan kunci ki. 

   Contoh: kunci = soreang       (dalam angka: 18, 14, 17, 4, 0, 13, 6)

 Plainteks: 

terimapesanangulaikambing

 Kunci: 

soreangsoreangsoreangsore

 Untuk 8 karakter berikutnya, kembali menggunakan pola enkripsi yang sama.

A = 0,  B = 1, C = 2,  D = 3,  E = 4,  F = 5,  G = 6,  H = 7,  I = 8,  J = 9,  K = 10

L = 11, M = 12, N = 13, O = 14, P = 15, Q = 16, R = 17, S = 18, T = 19, U = 20

V = 21, W = 22, X = 23, Y = 24, Z = 25

Ini artinya setiap huruf plainteks dienkripsi 

menggunakan Caesar Cipher dengan kunci k 

yang berbeda-beda

35

\end{document}
